% !TeX root = ../sections/02_exercises.tex
\subsection{2-A-4.}
\begin{exercise}
	\textup{(区間縮小法)}
	閉区間の列
	\[
	\{I_n=[a_n,b_n]\}_{n\in\mathbb{N}}
	\]
	が
	\[
	I_{n+1}\subset I_n
	\qquad
	(n=1,2,3,\ldots)
	\]
	をみたすとき，
	\[
	I:=\bigcap_{n=1}^{\infty} I_n
	\]
	は空集合ではないことを証明せよ。
	
	さらに，
	\[
	\lim_{n\to\infty}(b_n-a_n)=0
	\]
	ならば，
	\[
	I=\{\alpha\}
	\]
	と書くことができ，
	\[
	\lim_{n\to\infty}a_n
	=
	\lim_{n\to\infty}b_n
	=
	\alpha
	\]
	であることを証明せよ。
\end{exercise}

\begin{background}
	閉区間の列
	\[
	I_n=[a_n,b_n]
	\]
	が
	\[
	I_{n+1}\subset I_n
	\]
	を満たすとき，区間がだんだん小さくなっている。
	この条件から，
	\[
	a_n\leq a_{n+1}\leq b_{n+1}\leq b_n
	\]
	が成り立つ。
	
	したがって，\(\{a_n\}\) は単調増加数列であり，
	\(\{b_n\}\) は単調減少数列である。
	また，常に
	\[
	a_n\leq b_n
	\]
	である。
	
	単調増加かつ上に有界な数列は収束し，
	単調減少かつ下に有界な数列も収束する。
	これを用いて，共通部分に含まれる点を構成する。
\end{background}

\begin{strategy}
	まず，
	\[
	I_{n+1}\subset I_n
	\]
	から
	\[
	a_n\leq a_{n+1}\leq b_{n+1}\leq b_n
	\]
	を導く。
	これにより，\(\{a_n\}\) は単調増加，
	\(\{b_n\}\) は単調減少である。
	
	さらに，任意の \(n,m\) に対して
	\[
	a_n\leq b_m
	\]
	が成り立つことを確認する。
	これにより，\(\{a_n\}\) は上に有界である。
	したがって
	\[
	\alpha=\lim_{n\to\infty}a_n
	\]
	が存在する。
	
	次に，各 \(n\) について
	\[
	a_n\leq \alpha\leq b_n
	\]
	を示す。
	これが示せれば
	\[
	\alpha\in [a_n,b_n]=I_n
	\]
	が任意の \(n\) で成り立つので，
	\[
	\alpha\in\bigcap_{n=1}^{\infty}I_n
	\]
	となる。
	よって \(I\neq\emptyset\) である。
	
	さらに，
	\[
	b_n-a_n\to 0
	\]
	を仮定する。
	もし \(\alpha,\beta\in I\) ならば，任意の \(n\) に対して
	\[
	\alpha,\beta\in [a_n,b_n]
	\]
	であるから
	\[
	|\alpha-\beta|\leq b_n-a_n
	\]
	である。
	右辺は \(0\) に収束するので，
	\[
	|\alpha-\beta|=0
	\]
	となり，\(\alpha=\beta\) である。
	したがって \(I\) はただ一つの点からなる。
	
	最後に，その一点を \(\alpha\) と書けば，
	\[
	a_n\leq \alpha\leq b_n
	\]
	かつ
	\[
	b_n-a_n\to 0
	\]
	であるから，はさみうちにより
	\[
	a_n\to\alpha,
	\qquad
	b_n\to\alpha
	\]
	が従う。
\end{strategy}